% SECTION: Regression
\section{Regression}

\divider

% SUBSECTION: Wiederholung: Lineare Regression
\subsection{Wiederholung: Lineare Regression}

% Lineare Regression
\begin{frame}{Lineare Regression}{}
	Wir wiederholen die wichtigsten Aussagen zur linearen Regression:

	\begin{itemize}
		\item Wir betrachten einen Datensatz bestehend aus $N$ Featurevektoren $\bx_n \in \R^m, n = 1, 2, \ldots, N$.
		\item Jeder Vektor $\bx_n$ ist mit einem zugehörigen Label $y_n \in \R, n = 1, 2, \ldots, N$ versehen.
		\item Das lineare Regressionsmodell hat die Form $$h(\bx; \btheta) = \theta_0 + \sum_{m=1}^M \theta_m x_m.$$
		\item Im Folgenden bietet es sich an, die Dummy-Komponente $x_0 \coloneqq 1$ einzuführen,
			sodass gilt $$h(\bx; \btheta) = \sum_{m=0}^M \theta_m x_m = \btheta^\top \bx.$$
	\end{itemize}
\end{frame}


% Design Matrix
\begin{frame}{Design Matrix}{}
	\begin{itemize}
		\item Wir fassen die Featurevektoren $\bx_n$ zeilenweise zu der Matrix
		\begin{equation*}
			\bX \coloneqq
			\begin{pmatrix}
				\horzbar 	& \bx_1 	& \horzbar 	\\
				\horzbar 	& \bx_2 	& \horzbar 	\\
				\vdots 	&\vdots 	& \vdots 		\\
				\horzbar 	& \bx_N 	& \horzbar
			\end{pmatrix} \in \R^{N \times (M + 1)}
		\end{equation*}
			zusammen. Die erste Spalte ist hierbei konstant 1 \textit{(wegen der Dummy-Komponente)}.
		\item Wir nennen diese Matrix \keyword{Design Matrix}.
		\item Analog sammeln wir alle Labels in einem Vektor $$\by \coloneqq \bigl( y_1, y_2, \ldots, y_N \bigr)^\top \in \R^N.$$
	\end{itemize}
\end{frame}


% Methode der kleinsten Quadrate
\begin{frame}{Methode der kleinsten Quadrate}{}
	\begin{itemize}
		\item Der Vektor $\btheta \in \R^{M+1}$ beinhaltet die Modellparameter.
		\item Diese möchten wir so wählen, dass unser Modell die Daten möglichst gut erklärt,
			d.\,h. die Abweichung der Regressionsfunktion zu den Trainingsdaten sollte möglichst klein sein.
	\end{itemize}

	\begin{mydef}[def:kleinste-quadrate-fehler]
		Mit den obigen Definitionen ist der \keyword{Least-Squares-Fehler} gegeben durch
		\begin{equation*}
			\JLS(\btheta) \coloneqq \frac{1}{2N} \Vert \bX \btheta - \by \Vert^2.
		\end{equation*}
	\end{mydef}

	\textbf{Bemerkung:} Der Vektor $\bX \btheta \in \R^N$ beinhaltet die Modellvorhersagen,
	wenn man dem Modell die Parameter $\btheta$ zugrunde legt.
	$\JLS$ misst also die mittlere quadratische Abweichung der Vorhersagen zu den wahren Werten.
\end{frame}


% Minimierung der Fehlerfunktion
\begin{frame}{Minimierung der Fehlerfunktion}{}
	Unser Ziel ist die Minimierung von $\JLS$. Hierzu schreiben $\JLS$ unter Ausnutzung der Transpositionsregeln
	$(\bA + \bB)^\top = \bA^\top + \bB^\top$ und $(\bA \bB)^\top = \bB^\top \bA^\top$ um:
	\begin{align*}
		\JLS(\btheta)
			&= \frac{1}{2N} \Vert \bX \btheta - \by \Vert^2
				= \frac{1}{2N} \bigl( \bX \btheta - \by \bigr)^\top \bigl( \bX \btheta - \by \bigr)
	\\[2mm]
			&= \frac{1}{2N} \bigl( (\bX \btheta)^\top - \by^\top \bigr) \bigl( \bX \btheta - \by \bigr)
	\\[2mm]
			&= \frac{1}{2N} \bigl( (\bX \btheta)^\top \bX \btheta - (\bX \btheta)^\top \by - \by^\top (\bX \btheta) +
				\by^\top \by \bigr)
	\\[2mm]
			&= \frac{1}{2N} \bigl( \btheta^\top \bX^\top \bX \btheta - 2 (\bX \btheta)^\top \by + \by^\top \by \bigr)
	\\[2mm]
			&= \frac{1}{2N} \bigl( \btheta^\top \bX^\top \bX \btheta - 2 \btheta^\top \bX^\top \by + \by^\top \by \bigr)
	\end{align*}
\end{frame}


% Minimierung der Fehlerfunktion (Forts.)
\begin{frame}{Minimierung der Fehlerfunktion (Forts.)}{}
	Um einen stationären Punkt zu finden, müssen wir $\JLS$ nach $\btheta$ ableiten und Null setzen.

	\begin{itemize}
		\item Es ist $\frac{\partial}{\partial \btheta} \btheta^\top \bX^\top \bX \btheta = 2 \bX^\top \bX \btheta$,
			da $\bX^\top \bX$ symmetrisch ist. \\ \vspace*{2mm}
		\begin{myproof}
			Wir nutzen die Transpositionsregeln und erhalten $\bigl( \bX^\top \bX \bigr)^\top =
			\bX^\top \bigl( \bX^\top \bigr)^\top = \bX^\top \bX$, woraus die Symmetrie folgt.
		\end{myproof} \vspace*{2mm}
		\item Es ist $\frac{\partial}{\partial \btheta} \bigl( -2 \btheta^\top \bX^\top \by \bigr) = -2 \bX^\top \by$
			und $\frac{\partial}{\partial \btheta} \by^\top \by = \0$.
		\item Der Gradient der Fehlerfunktion $\JLS$ ist damit gegeben durch
		\begin{align*}
			\frac{\partial}{\partial \btheta} \JLS(\btheta)
				&= \frac{1}{2 N} \bigl( 2 \bX^\top \bX \btheta - 2 \bX^\top \by \bigr)
					\propto \bX^\top \bX \btheta - \bX^\top \by
		\\[2mm]
				&= \bX^\top \bigl( \bX \btheta - \by \bigr)
		\end{align*}
	\end{itemize}
	\textbf{Bemerkung:} Man beachte den Zusammenhang zwischen dem Gradienten und dem Vorhersagefehler.
\end{frame}


% Minimierung der Fehlerfunktion (Forts.)
\begin{frame}{Minimierung der Fehlerfunktion (Forts.)}{}
	Wir setzen den Gradienten nun Null und lösen nach $\btheta$ auf:
	\begin{alignat*}{3}
		\frac{\partial}{\partial \btheta} \JLS(\btheta) \overset{!}{=} \0
			&\quad\Longleftrightarrow\quad
			&\bX^\top \bX \btheta - \bX^\top \by
			&= \0
	\\[2mm]
			&\quad\Longleftrightarrow\quad
			&\bX^\top \bX \btheta
			&= \bX^\top \by
	\\[2mm]
			&\quad\Longleftrightarrow\quad
			&\btheta^\star
			&= \bigl( \bX^\top \bX \bigr)^{-1} \bX^\top \by
	\end{alignat*}

	\textbf{Bemerkung:} Wir haben damit einen stationären Punkt von $\JLS$ gefunden. Allerdings wissen wir noch nicht,
	ob es sich tatsächlich um ein Minumum handelt. Außerdem sind wir hier etwas naiv von der Existenz der Inversen
	$\bigl( \bX^\top \bX \bigr)^{-1}$ ausgegangen. Existiert diese Matrix immer? \\
	\textit{Mit diesen Fragen beschäftigen wir uns im Folgenden.}
\end{frame}


% Ein paar nützliche Aussagen
\begin{frame}{Ein paar nützliche Aussagen}{}
	\begin{mytheo}
		Für alle $\bA \in \R^{N \times M}$ ist $\bA^\top \bA$ positiv semidefinit.
	\end{mytheo}
	
	\begin{myproof}
		Für alle $\bz \in \R^M$ gilt
		$\bz^\top \bigl( \bA^\top \bA \bigr) \bz = (\bA \bz)^\top (\bA \bz) = \Vert \bA \bz \Vert^2 \ge 0$,
		die Behauptung.	
	\end{myproof}

	\begin{mytheo}[th:spaltenrang-pd]
		Für alle $\bA \in \R^{N \times M}$ mit \textit{vollem Spaltenrang} ist $\bA^\top \bA$ positiv definit.
	\end{mytheo}

	\begin{myproof}
		Da $\bA$ vollen Spaltenrang hat, sind die Spaltenvektoren linear unabhängig.
		Daher ist $\bA \bz$ für alle $\bz \in \R^M \backslash \{ \0 \}$ eine nicht-triviale Linearkombination
		der Spalten von $\bA$. Folglich gilt $\bA \bz \ne \0$ und wir erhalten
		$\bz^\top \bigl( \bA^\top \bA \bigr) \bz = (\bA \bz)^\top (\bA \bz) = \Vert \bA \bz \Vert^2 > 0.$
	\end{myproof}
\end{frame}


% Ein paar nützliche Aussagen (Forts.)
\begin{frame}{Ein paar nützliche Aussagen (Forts.)}{}
	\begin{mytheo}[th:spaltenrang-inv]
		Wenn $\bA$ positiv definit ist, so ist $\bA$ auch invertierbar.
	\end{mytheo}
	
	\begin{myproof}
		Wir nehmen an $\bA$ sei nicht invertierbar. Dann sind die Spalten von $\bA$ linear abhängig und es existiert ein
		$\bz \in \R^M \backslash \{ \0 \}$ mit $\bA \bz = \0$.
		Daraus folgt aber $$\bz^\top \bA \bz = \bz^\top \0 = 0,$$ was der Definitheitseigenschaft von $\bA$ widerspricht.
	\end{myproof}
\end{frame}


% Handelt es sich nun um ein Minimum?
\begin{frame}{Handelt es sich nun um ein Minimum?}{}
	Um diese Frage zu klären, berechnen wir die Hesse-Matrix von $\JLS$ und untersuchen,
	ob die Matrix im stationären Punkt positiv definit ist.

	\begin{itemize}
		\item Man rechnet nach, dass die Hesse-Matrix gegeben ist durch
			$$\frac{\partial^2}{\partial \btheta \partial \btheta} \JLS(\btheta) = \bX^\top \bX.$$
		\item Mit Satz \ref{th:spaltenrang-pd} ist diese Matrix positiv definit für alle $\btheta$, falls $\bX$ vollen Spaltenrang hat,
			d.\,h. die Features müssen linear unabhängig sein.
			In diesem Fall handelt es sich bei $\btheta^\star$ also tatsächlich um ein Minimum.
		\item Unter derselben Voraussetzung haben wir damit sogar die \textbf{Konvexität} von $\JLS$ bewiesen,
			sodass das Mimimun global und eindeutig ist.
	\end{itemize}
\end{frame}


% Und wie sieht es nun mit der Existenz der Inversen aus?
\begin{frame}{Und wie sieht es nun mit der Existenz der Inversen aus?}{}
	\begin{itemize}
		\item Mit Satz \ref{th:spaltenrang-inv} ist die Existenz der Inversen sichergestellt, wenn $\bX^\top \bX$ positiv definit ist,
			was wiederum aus der linearen Unabhängigkeit der Features folgt (Satz \ref{th:spaltenrang-pd}).
		\item \textbf{Achtung:} Bei der numerischen Berechnung der Inversen auf dem Computer kann es jedoch
			trotzdem zu Problemen kommen. Es ist ratsam, die Inverse nicht explizit zu berechnen,
			sondern stattdessen das lineare Gleichungssystem $\bX^\top \bX \btheta = \bX^\top \by$ nach $\btheta$ zu lösen.
			Dieses Vorgehen ist numerisch stabiler.
		\item Alternativ kann das Optimierungsproblem mit dem \keyword{Gradientenabstiegsverfahren} gelöst werden.
			Iteriere hierzu die Regel:
		\begin{equation}
			\btheta_{t+1} \longleftarrow \btheta_t - \alpha \bX^\top \bigl( \bX \btheta_t - \by \bigr). 
		\end{equation}
	\end{itemize}
\end{frame}


% Normalengleichungen
\begin{frame}{Normalengleichungen}{}
	Wir fassen unsere Ergebnisse zusammen:
	
	\begin{mytheo}[th:normal-equations]
		Die Design Matrix $\bX \in \R^{N \times (M+1)}$ des Regressionsproblems habe vollen Spaltenrang,
		d.\,h. die Features seien linear unabhängig.
		Dann existiert die Inverse $\bigl( \bX^\top \bX \bigr)^{-1}$ und
		$$\btheta^\star = \bigl( \bX^\top \bX \bigr)^{-1} \bX^\top \by$$ ist das eindeutige globale Minimum von $\JLS$.
		Die Gleichungen $\bX^\top \bX \btheta = \bX^\top \by$ bezeichnet man als \keyword{Normalengleichungen}.
	\end{mytheo}
\end{frame}


% Probabilistische Regression
\subsection{Probabilistische Regression}

% Eine probabilistische Sicht
\begin{frame}{Eine probabilistische Sicht}{}
	\begin{boxGray}
		\textbf{Annahme 1:} Die Zielwerte $y \in \R$ und die Eingaben $\bx \in \R^{m+1}$ hängen über die Gleichung
		\begin{equation}
			y = \btheta^\top \bx + \varepsilon
		\end{equation}
		zusammen. Dabei ist $\varepsilon \in \R$ ein Fehlerterm, der nicht modellierte Effekte oder Datenrauschen repräsentiert.
	\end{boxGray}
	
	\begin{boxGray}
		\textbf{Annahme 2:} Für den Fehlerterm gilt $\varepsilon \sim \norm\bigl( 0, \sigma^2 \bigr)$.
	\end{boxGray}

	Der Zielwert $y$ kann damit als Zufallsvariable aufgefasst werden, für die gilt
	\begin{equation}
		y \mid \bx \sim \norm\bigl(\btheta^\top \bx, \sigma^2 \bigr).
	\end{equation}
\end{frame}


% Visualisierung: Probabilistische Regression
\begin{frame}{Visualisierung: Probabilistische Regression}{}
	tbd
\end{frame}


% Maximum-Likelihood Regression
\begin{frame}{Maximum-Likelihood Regression}{}
	\begin{itemize}
		\item Wir betrachten wieder einen Datensatz $\bigl\{ (\bx_n, y_n) \bigr\}$ bestehend aus $N$ Trainingsbeispielen.
		\item Die Likelihood des Modells ist gegeben durch
		\begin{equation}
			\label{eq:reg-likelihood}
			\begin{aligned}
				p\bigl( \by \mid \bX; \btheta, \sigma^2 \bigr)
					&= \prod_{n=1}^N p\bigl( y_n \mid \bx_n; \btheta, \sigma^2 \bigr)
			\\[2mm]
					&= \prod_{n=1}^N \frac{1}{\sqrt{2 \pi \sigma^2}} \cdot \exp\left(
						-\frac{1}{2 \sigma^2} \biggl( y_n - \btheta^\top \bphi(\bx_n) \biggr)^2
					\right)
			\end{aligned}
		\end{equation}
	\end{itemize}
\end{frame}


% Maximum-Likelihood Regression (Forts.)
\begin{frame}{Maximum-Likelihood Regression (Forts.)}{}
	\begin{itemize}
		\item Durch Logarithmieren von \eqref{eq:reg-likelihood} erhalten wir die Log-Likelihood Funktion
		\begin{equation}
			\label{eq:reg-log-likelihood}
			\begin{aligned}
				l\bigl( \btheta, \sigma^2 \bigr)
					&\coloneqq \ln p\bigl( \by \mid \bX; \btheta, \sigma^2 \bigr)
			\\[2mm]
					&= -\frac{N}{2} \ln\bigl( 2 \pi \sigma^2 \bigr)
						- \frac{1}{2 \sigma^2} \sum_{n=1}^N \bigl( y_n - \btheta^\top \bphi(\bx_n) \bigr)^2.
			\end{aligned}
		\end{equation}
		\item Aus \eqref{eq:reg-log-likelihood} wird ersichtlich, dass der Least-Squares-Fehler minimiert werden muss, um die
			Likelihood in $\btheta$ zu maximieren.
		\item Die Lösung für $\btheta$ ist dieselbe wie in Satz \ref{th:normal-equations}, zusätzlich kann die globale
			Unsicherheit des Modells quantifiziert werden
		\vspace*{-3mm}
		\begin{equation}
			\bthetaestML = \bigl( \bPhi^\top \bPhi \bigr)^{-1} \bPhi^\top \by
		\qquad\qquad
			\sigmasquML = \frac{1}{N} \sum_{n=1}^N \bigl( y_n - \bthetaestML^\top \bphi(\bx_n) \bigr)^2.
		\end{equation}
	\end{itemize}
\end{frame}


% Visualisierung: Maximum-Likelihood Regression
\begin{frame}{Visualisierung: Maximum-Likelihood Regression}{}
	tbd
\end{frame}


% Probabilistische Regression
\subsection{Maximum Aposteriori (MAP) Regression}

% Berücksichtigung einer Apriori-Verteilung
\begin{frame}{Berücksichtigung einer Apriori-Verteilung}{}
	\begin{itemize}
		\item Die Maximum-Likelihood Lösung ist \textbf{anfällig für Overfitting}.
		\item Eine Möglichkeit der \keyword{Regularisierung} ist die Verwendung einer \keyword{Apriori-Verteilung} (Prior).
		\item Für diese Verteilung benutzen wir eine Normalverteilung für die Parameter $\btheta$
		\begin{equation}
			\btheta \sim \norm\bigl( \0, b^2 \bI \bigr).
		\end{equation}
		\item Es handelt sich hierbei um eine \keyword{isotropische Normalverteilung} (rotationsinvariant).
		\item Der Prior kann benutzt werden, um Vorwissen zu berücksichtigen.
		\item Er kodiert, welche Parameterwerte plausibel sind.
		\item Häufig ist der Prior jedoch \keyword{nicht informativ} (z.\,B. Gleichverteilung).
	\end{itemize}
\end{frame}


% Warum die Normalverteilung?
\begin{frame}{Warum die Normalverteilung?}{}
	\begin{boxGray}
		\textbf{Frage:} Warum benutzen wir eine Normalverteilung?
	\end{boxGray}
	
	\begin{itemize}
		\item Ein Argument ist \keyword{Konjugiertheit}: Da Prior und Likelihood durch Normalverteilungen gegeben sind, ist auch der \textit{Posterior normalverteilt}.
		\item Für Operationen auf Normalverteilungen gibt es meistens \textit{analytische Lösungen}.
	\end{itemize}
\end{frame}


% Satz von Bayes
\begin{frame}{Satz von Bayes}{}
	\begin{itemize}
		\item Wir benutzen den \keyword{Satz von Bayes}
		\begin{equation} \label{eq:reg-bayes-theorem}
			\begin{aligned}
				p\bigl( \btheta \mid \bX, \by \bigr)
					&= \frac{\overbracket{
						p\bigl( \by \mid \bX, \btheta; \sigma^2 \bigr)
					}^{\text{Likelihood}} \cdot \overbracket{p\bigl( \btheta; b^2 \bigr)}^{\text{Prior}}}{p\bigl( \by \mid \bX \bigr)}
			\\[2mm]
					&\propto p\bigl( \by \mid \bX, \btheta; \sigma^2 \bigr) \cdot p\bigl( \btheta; b^2 \bigr).
			\end{aligned}
		\end{equation}
		\item Wir betrachten $\btheta$ nun als (vektorwertige) Zufallsvariable.
		\item Bisher haben wir nur die Likelihood maximiert, nun berücksichtigen wir auch den Prior $p\bigl( \btheta; b^2 \bigr)$.
	\end{itemize}
	
	\begin{boxGray}
		\textbf{Frage:} Warum ignorieren wir den Nenner in \eqref{eq:reg-bayes-theorem}?
	\end{boxGray}
\end{frame}


% Log-Posterior Verteilung
\begin{frame}{Log-Posterior Verteilung}{}
	\begin{itemize}
		\item Anstelle der Log-Likelihood maximieren wir den \keyword{Log-Posterior}.
		\item Diesen erhalten wir aus \eqref{eq:reg-bayes-theorem}
		\begin{equation}
			\ln p\bigl( \btheta \mid \bX, \by \bigr)
				= \ln p\bigl( \by \mid \bX, \btheta; \sigma^2 \bigr) + \ln p\bigl( \btheta; b^2 \bigr) + \const.
		\end{equation}
		\item Die Konstante $\const$ enthält alle Terme, die nicht von $\btheta$ abhängen.
		\item Die \keyword{MAP-Schätzung} (Maximum Aposteriori) ist ein Kompromiss zwischen dem Prior und der datenabhängigen Likelihood.
		\item Wir betrachten im Folgenden den \keyword{negativen Log-Posterior} und erhalten
		\begin{equation} \label{eq:reg-neg-log-posterior}
			\bthetaestMAP \coloneqq \argmin_{\btheta}\Bigl\{ -\ln p\bigl( \by \mid \bX, \btheta; \sigma^2 \bigr) - \ln p\bigl( \btheta; b^2 \bigr) \Bigr\}.
		\end{equation}
	\end{itemize}
\end{frame}


% Log-Posterior Verteilung (Forts.)
\begin{frame}{Log-Posterior Verteilung (Forts.)}{}
	Unter Verwendung der Transpositionsregeln schreiben wir \eqref{eq:reg-neg-log-posterior} um zu
	\begin{equation} \label{eq:reg-neg-log-posterior-neu}
		\begin{aligned}
			-\ln p\bigl( \btheta \mid \bX, \by \bigr)
				&= -\ln p\bigl( \by \mid \bX, \btheta; \sigma^2 \bigr) - \ln p\bigl( \btheta; b^2 \bigr) + \const
		\\[2mm]
				&= \frac{1}{2 \sigma^2} \bigl( \by - \bPhi \btheta \bigr)^\top \bigl( \by - \bPhi \btheta \bigr) + \frac{1}{2 b^2} \btheta^\top \btheta + \const
		\\[2mm]
				&= \frac{1}{2 \sigma^2} \bigl( \by^\top \by - \by^\top \bPhi \btheta - (\bPhi \btheta)^\top \by + (\bPhi \btheta)^\top \bPhi \btheta \bigr)
					+ \frac{1}{2 b^2} \btheta^\top \btheta + \const
		\\[2mm]
				&= \frac{1}{2 \sigma^2} \bigl( \btheta^\top \bPhi^\top \bPhi \btheta - 2 \btheta^\top \bPhi^\top \by \bigr)
					+ \frac{1}{2 b^2} \btheta^\top \btheta + \const.
		\end{aligned}
	\end{equation}
	
	\begin{mytask}
		Überlegen Sie sich, wie man in \eqref{eq:reg-neg-log-posterior-neu} von Zeile 1 auf Zeile 2 kommt.
	\end{mytask}
\end{frame}


% Berechnung des Gradienten
\begin{frame}{Berechnung des Gradienten}{}
	\begin{mytask}
		Jetzt sind Sie an der Reihe:
		\begin{enumerate}
			\item Berechnen Sie den Gradienten der negativen Log-Likelihood $-\ln p\bigl( \btheta \mid \bX, \by \bigr)$.
				Benutzen Sie hierzu die Form aus \eqref{eq:reg-neg-log-posterior-neu}.
			\item Setzen Sie den Gradienten anschließend Null und lösen Sie nach $\btheta$.
			\item Die Lösungsformel für $\bthetaestMAP$ fordert die Invertierung einer Matrix. Zeigen Sie, dass diese Inverse (zumindest analytisch) immer existiert.
		\end{enumerate}
	\end{mytask}
\end{frame}


% Vergleich der MLE und MAP Lösungen
\begin{frame}{Vergleich der MLE und MAP Lösungen}{}
	\begin{boxGray}
		\textbf{MLE Schätzer:}
		\begin{equation}
			\bthetaestML = \bigl( \bPhi^\top \bPhi \bigr)^{-1} \bPhi^\top \by
		\end{equation}
	\end{boxGray}
	
	\begin{boxGray}
		\textbf{MAP Schätzer:}
		\begin{equation}
			\bthetaestMAP = \left( \bPhi^\top \bPhi + \frac{\sigma^2}{b^2} \bI \right)^{-1} \bPhi^\top \by
		\end{equation}
	\end{boxGray}
\end{frame}


% Visualisierung: Vergleich der MLE und MAP Lösungen
\begin{frame}{Visualisierung: Vergleich der MLE und MAP Lösungen}{}
	tbd
\end{frame}














